\section*{Preface}

This Master's thesis was conducted at the Institute of Industrial Automation and Software Engineering (IAS) at the University of Stuttgart. The research presented in this work investigates the intersection of embodied artificial intelligence, edge computing, and robotics perception, focusing on how inference-time visual prompt engineering and dynamic video keyframe selection can enhance scene comprehension in Vision-Language Models (VLMs) deployed on mobile robotics platforms.

The primary motivation behind this work stems from a practical challenge in autonomous systems: while multimodal foundation models exhibit exceptional reasoning capabilities, their deployment on resource and energy-constrained edge devices remains bound by strict latency, memory, and thermal limits. By exploring non-invasive visual transformations and container-level video extraction pipelines, this research aims to bridge the gap between heavy cloud-level reasoning architectures and real-time onboard robotic control.

\section*{Acknowledgments}

First and foremost, I would like to express my sincere gratitude to \textbf{Prof. Dr.-Ing. Andrey Morozov} for granting me the opportunity to conduct this research at the Institute of Industrial Automation and Software Engineering (IAS) and for providing the academic vision, resources, and environment essential for completing this work. 

I am deeply grateful to my supervisor, \textbf{Joachim Grimstad}, for his continuous guidance, insightful discussions, and dedicated support throughout this project. His constructive feedback, technical expertise, and encouraging advice were instrumental in refining the methodology, experimental rigor, and presentation of this thesis.

I also extend my appreciation to the researchers, colleagues, and technical staff at the University of Stuttgart for maintaining the high-performance computing clusters and edge robotics hardware platforms that made the extensive empirical benchmarks and multi-agent evaluations possible.

